% Derivation of the two-level unified model (route 2 of P3_OUTLINE.ja.md).
% Build: pdflatex unified_model_derivation.tex (twice).
% Sources: Klempt et al. 2024 (references/, Eq. 20-36) and Klempt et al. 2026 (the point model,
% as coded in JAXFEM/core_hamilton_1d.py). Nothing from F. Klempt's unpublished work.
\documentclass[11pt,a4paper]{article}
\usepackage[margin=2.2cm]{geometry}
\usepackage{amsmath,amssymb,amsthm}
\usepackage{mathptmx}
\usepackage{booktabs}
\newtheorem{proposition}{Proposition}
\newtheorem{remark}{Remark}
\newcommand{\dd}{\,\mathrm{d}}
\begin{document}
\begin{center}
{\large\bfseries A two-level biofilm model from one extended Hamilton principle}\\[3pt]
Keisuke Nishioka, working note, 15 October 2026 (not part of the thesis)
\end{center}

\paragraph{Aim.} In the thesis the growth field of Klempt et al.\ 2024 and the point
model of Klempt et al.\ 2026 are solved separately and joined at the Gauss points.
This needs choices that no paper fixes: how the void of the point model relates to the
void of the field, whether the point model is scaled to the field, what drives the
growth, and how the two clocks are linked. Here all equations follow from one energy,
one dissipation potential and one constraint.

\section{Why one level is not enough}

The direct generalisation takes every species as a field $\phi_i(\mathbf X,t)$ with one
void $\phi_0=1-\sum_i\phi_i$ and the point-model energy at every point. The point model
keeps every $\phi_i$ away from zero by its barrier energy
$K/[\phi_i^2(1-\phi_i)^2]$,
and its production is local. A 1D solution with all fields in one Newton system
shows the consequence: the outer liquid
fills up from the barrier floor everywhere at once, so the model cannot represent a
biofilm next to liquid. The two voids are therefore different physical objects: the
liquid outside the biofilm, and the liquid inside it. This motivates two levels.

\section{State, energy, dissipation}

\begin{tabular}{@{}ll@{}}
$\phi(\mathbf X,t)\in[0,1]$ & biofilm phase (Klempt et al.\ 2024): biofilm against outer liquid\\
$s_i,\ s_0$ & composition of the biofilm: species $i=1\dots n$ and inner liquid, $\sum_i s_i+s_0=1$\\
$\psi_i$ & living fraction of species $i$ (Klempt et al.\ 2026)\\
$\alpha$ & local expansion, $\mathbf F_g=\alpha\mathbf I$, $\alpha(0)=1$\\
$\phi s_i$, $\phi\sum_i s_i\psi_i$ & absolute amount of species $i$, living biomass
\end{tabular}

\medskip
With $w=\phi+\varepsilon$ ($\varepsilon>0$ small, it keeps the composition equations
regular in the outer liquid where they carry no weight), the energy density is
\begin{equation}
E=\frac{\beta}{2}|\nabla\phi|^2+\lambda w\,\Psi_{\rm pm}(s,s_0,\psi)
+w\frac{\beta_s}{2}\sum_i|\nabla s_i|^2-k_\alpha(\alpha-1)\,\phi\sum_i s_i\psi_i
+E_{\rm mech},
\label{eq:E}
\end{equation}
with the point-model energy of Klempt et al.\ 2026
\begin{equation}
\Psi_{\rm pm}=K\Big[\sum_iB(s_i)+B(s_0)+\sum_iB(\psi_i)\Big]
+\frac{\alpha^*}{2}\sum_ib_i\psi_i^2-\frac{c^*}{2}\sum_{ij}A_{ij}\,s_i\psi_i\,s_j\psi_j,
\qquad B(x)=\frac{1}{x^2(1-x)^2},
\end{equation}
and $E_{\rm mech}=\tfrac{\mu}{2}\phi^2(\mathbf I:\mathbf C_e-3)$ as in Klempt et al.\
2024 (dropped below for the same reason as there: its effect on growth is small). The
dissipation potential is
\begin{equation}
D=\frac{\eta_\phi}{2}\dot\phi^2+\lambda w\,D_{\rm pm}+\frac{\eta_\alpha}{2}\dot\alpha^2,\qquad
D_{\rm pm}=\sum_i\Big[\frac{\eta}{2}\dot s_i^2+\frac{\eta}{2}\big(\tfrac{\dd}{\dd t}(s_i\psi_i)\big)^2\Big]+\frac{1}{2}\dot s_0^2,
\label{eq:D}
\end{equation}
and the constraint is imposed with the multiplier field $\gamma$ through
$\lambda w\gamma(\sum_is_i+s_0-1)$. The only new parameters are $\lambda$ (the energy
scale of the composition relative to the phase), $\beta_s$ (spreading of the composition
inside the biofilm) and the form of the growth energy.

\section{Stationarity conditions}

The extended Hamilton principle
\[\delta_q\int_\Omega E\dd V+\int_\Omega\partial_{\dot q}D\,\delta q\dd V=0\]
for all variations $\delta q$ of $q=\{\phi,s_i,s_0,\psi_i,\alpha,\gamma\}$, with
$\mathbf n\cdot\nabla\phi=\mathbf n\cdot\nabla s_i=0$ on the boundary, gives:

\paragraph{Composition} (divided by $\lambda w>0$):
\begin{align}
K B'(s_i)+\gamma+\eta(1+\psi_i^2)\dot s_i+\eta s_i\psi_i\dot\psi_i
-c^*\psi_i\Big(\sum_jA_{ij}s_j\psi_j+\frac{k_\alpha(\alpha-1)\phi}{\lambda w\,c^*}\Big)
-\frac{\beta_s}{\lambda w}\nabla\!\cdot(w\nabla s_i)&=0,\\
K B'(s_0)+\gamma+\dot s_0&=0,\\
K B'(\psi_i)+\alpha^*b_i\psi_i+\eta s_i\psi_i\dot s_i+\eta s_i^2\dot\psi_i
-c^*s_i\Big(\sum_jA_{ij}s_j\psi_j+\frac{k_\alpha(\alpha-1)\phi}{\lambda w\,c^*}\Big)&=0,\\
\textstyle\sum_is_i+s_0-1&=0.
\end{align}
Without the growth term and without composition gradients these are exactly the point
model of Klempt et al.\ 2026 at every point. The growth energy enters with the same
structure as the interactions.

\paragraph{Phase:}
\begin{equation}
\eta_\phi\dot\phi-\beta\nabla^2\phi+\lambda\Psi_{\rm pm}+\frac{\beta_s}{2}\sum_i|\nabla s_i|^2
-k_\alpha(\alpha-1)\sum_is_i\psi_i=0 .
\end{equation}
The composition acts on the phase through $\lambda\Psi_{\rm pm}$: where the composition
energy is low (strong favourable interactions), the biofilm tends to grow, and the growth
source is weighted by the living share. The dependence of the dissipation weight $w$ on
$\phi$ is not a force: the principle varies $D$ with respect to the rates only.

\paragraph{Expansion:}
\begin{equation}
\eta_\alpha\dot\alpha=k_\alpha\,\phi\sum_is_i\psi_i .
\label{eq:alpha}
\end{equation}

\section{Limits}

\begin{enumerate}
\item $\lambda\to0$, $\beta_s=0$: the composition at every point is the point model on
  its own clock, started where the biofilm is, and \eqref{eq:alpha} is the growth law of
  the element's scheme (c), $\dot\alpha=\kappa\,\phi_{3D}\sum_i\bar\phi_i$ with
  $\kappa=k_\alpha/\eta_\alpha$. Scheme (c) is thus the decoupled limit of this model,
  and its point model, which runs unscaled as in scheme (b), is the consistent one: its
  $\phi_0$ is the inner liquid $s_0$.
\item With $\phi\sum_is_i$ in place of the living biomass in the growth energy and one
  species, \eqref{eq:alpha} is Eq.~36 of Klempt et al.\ 2024 up to the factor $\sum_is_i$.
\item Scheme (a), which scales the point model to $\sum_i\phi_i=\phi_{3D}$, identifies
  the inner liquid with the outer one ($s_0=1-\phi$). It is not a limit of the model.
\end{enumerate}

\section{Energy}

\begin{proposition}
Let $E$ not depend on time explicitly and let no external power act. Along solutions,
\[\frac{\dd}{\dd t}\int_\Omega E\dd V=-2\int_\Omega D\dd V\le0 .\]
\end{proposition}
\begin{proof}
Every stationarity condition has the form $\partial_qE+(\text{gradient terms})
+\partial_{\dot q}D+(\text{constraint terms})=0$. Multiply each by $\dot q$, integrate
over $\Omega$ and add. The gradient terms integrate by parts to the rate of the gradient
energies (the boundary terms vanish by the no-flux conditions); the constraint terms give
$\lambda w\gamma\,\frac{\dd}{\dd t}(\sum_is_i+s_0)=0$; and $D$ is quadratic in the rates,
so $\sum_q\dot q\,\partial_{\dot q}D=2D$. What remains is
$\frac{\dd}{\dd t}\int E=-\int2D$.
\end{proof}

\begin{remark}
Switching an antibiotic on ($\alpha^*$ changing in time) adds
$\int\partial_tE=\int\lambda w\tfrac{\dot\alpha^*}{2}\sum_ib_i\psi_i^2$, an external
input; this is the jump seen at the switch in the 1D runs. Klempt et al.\ 2024's front
term $h^*_\phi$ is not derived from a potential, so it is left out here; a model that
includes it is no longer a gradient flow.
\end{remark}

\section{Discretisation and checks}

1D, nodes with mirrored ends, all unknowns of all nodes in one residual, implicit Euler,
full Newton to $|R|<10^{-10}$ with damping that keeps $s,\psi\in(0,1)$ and $\phi>0$
(\texttt{JAXFEM/unified\_two\_level\_1d.py}). For the one-level model the same scheme
reproduces the converged point model to $1.6\times10^{-13}$ and the discrete energy
decreases at every step (\texttt{unified\_1d\_monolithic.py}).

\section{What remains assumed}

The form of the growth energy (living biomass, and $\alpha-1$ so that it vanishes before
growth); $\lambda$ and $\beta_s$; the clock of the composition relative to the phase,
which is now the ratio of the dissipation coefficients $\eta/\eta_\phi$ and can be fixed
by the measured volume; the stiffness; detachment.
\end{document}
